\documentclass{article}
\usepackage{iclr2026_conference,times}   % swap to the ICLR 2027 style when released

\input{math_commands.tex}

\usepackage{hyperref}
\usepackage{url}
\usepackage{amsmath}
\usepackage{amssymb}
\usepackage{booktabs}
\usepackage{graphicx}
\usepackage{amsthm}
\usepackage{pifont}

\newcommand{\cmark}{\ding{51}}
\newcommand{\xmark}{\ding{55}}
\newcommand{\pmark}{\ensuremath{\sim}}

\newtheorem{proposition}{Proposition}
\newtheorem{assumption}{Assumption}

\title{PRISM: Structurally Reversible Latent Transport with \\ Feasibility by Construction for Ill-Posed Inverse Problems}

\author{Anonymous authors \\ Paper under double-blind review}

\begin{document}

\maketitle

% ================= Abstract (Phase 0 reframed) =================
\begin{abstract}
Recovering physical parameters from indirect, noisy observations---the ill-posed
inverse problem---requires properties that contemporary learning-based methods
provide only in isolation: consistent inversion, satisfaction of hard physical
constraints, and calibrated posterior uncertainty. We introduce PRISM, a
continuous normalizing-flow framework that combines structurally reversible latent
transport, explicit support-aware parameter transforms, and amortized conditional
posterior inference within a single model. The transport is a liquid
time-constant neural ordinary differential equation (ODE): under stated regularity
assumptions its time-$1$ map is a $C^{1}$ diffeomorphism, so the inverse is
obtained by reverse-time integration and forward--reverse reconstruction is
limited only by the numerical integration scheme. Fixed, dimension-matched support
transforms guarantee that every decoded parameter is physically feasible by
construction, with correct physical-space likelihoods, and an observation-conditioned
base distribution amortizes uncertainty-aware inference in a single forward pass.
We analyze the reversibility and stability of the latent transport under explicit
regularity conditions, and evaluate PRISM on partial-differential-equation (PDE)
operator inversion, simulation-based inference (SBI), and dynamical-system
benchmarks against strong, identically constrained baselines with multi-seed
uncertainty estimates. PRISM delivers strong inverse accuracy and competitive
posterior calibration under a shared training framework that also provides
reversible latent transport and feasibility by construction. We release a
reproducible benchmark and a scikit-learn-style library.
\end{abstract}

% ================= Introduction =================
\section{Introduction}
Inverse problems pervade the computational sciences: reconstructing a subsurface
field from sparse gauges, inferring dynamical parameters from a short trajectory,
or recovering the latent drivers of a stochastic simulator. They are ill-posed---the
forward operator destroys information---and their solutions are rarely
unconstrained. A useful method must therefore invert consistently, respect the
constraints that define an admissible solution, and quantify uncertainty honestly.

Existing families each address a strict subset. Neural operators and physics-informed
networks learn accurate forward surrogates but provide neither an inverse map nor a
posterior; invertible networks and flows furnish an exact inverse and a sampleable
posterior but enforce constraints only through soft penalties; and neural posterior
estimators target calibration directly yet inherit neither consistency nor
feasibility guarantees. We show these properties can be obtained together through
architectural choices rather than additional penalties.

We are careful to separate what is \emph{guaranteed by construction} from what is
\emph{demonstrated empirically}. PRISM guarantees reversibility of the latent
transport (under stated regularity) and feasibility through explicit support
transforms; it demonstrates calibration, inverse accuracy, and robustness on
benchmarks. It does not claim exactness between the auxiliary physical
forward surrogate and the inverse path---a quantity we instead measure.

\paragraph{Contributions.}
\begin{itemize}
\item A liquid neural-ODE transport for reversible latent parameter mapping, with
reversibility and stability analyzed under explicit regularity assumptions.
\item Explicit, dimension-matched support transforms that make posterior samples
physically feasible by construction, with the corresponding physical-space
likelihood.
\item An observation-conditioned base distribution for amortized posterior
inference in a single forward pass.
\item A cross-domain empirical study against \emph{identically constrained}
baselines with multi-seed uncertainty, separating guaranteed from empirical
properties.
\item A reproducible benchmark and documented library that regenerate every result.
\end{itemize}

% ================= Related Work =================
\section{Related Work}
\label{sec:related}
Fourier Neural Operators (FNOs), Deep Operator Networks (DeepONet), and the neural
operator framework learn resolution-invariant forward maps; physics-informed
neural networks (PINNs) embed PDE residuals. None inverts or yields a posterior.
Normalizing flows and continuous normalizing flows give tractable densities and
exact inversion; conditional invertible neural networks (cINNs) and related
amortized methods recover posteriors by sampling a latent code, but rarely enforce
hard constraints and are exactly consistent only for discrete coupling
architectures. Neural posterior estimation (NPE) targets calibrated posteriors,
assessed by simulation-based calibration (SBC), classifier two-sample tests
(C2ST), and standardized suites. Liquid time-constant networks introduce
input-dependent time constants; we repurpose this mechanism as an invertible flow
and isolate its contribution.

\begin{table}[t]
\centering
\caption{Architectural capability matrix. Each entry follows from the method
definition. \cmark{}~present, \xmark{}~absent by default. A support transform
(row~3) can be added to any method; we do so for all baselines in the constraint
comparison, so feasibility is a shared, testable property rather than unique to
PRISM.}
\label{tab:novelty}
\begin{tabular}{lcccccc}
\toprule
Capability & PRISM & cINN & NPE-CNF & VI & FNO/DeepONet & PINN \\
\midrule
Invertible latent transform      & \cmark & \cmark & \cmark & \cmark & \xmark & \xmark \\
Tractable conditional density    & \cmark & \cmark & \cmark & \cmark & \xmark & \xmark \\
Explicit support transform       & \cmark & \xmark & \xmark & \xmark & \xmark & \xmark \\
Auxiliary physical forward map    & \cmark & \xmark & \xmark & \xmark & \cmark & \cmark \\
Amortized single-pass sampling   & \cmark & \cmark & \cmark & \cmark & \xmark & \xmark \\
Continuous-time (ODE) transport  & \cmark & \xmark & \cmark & \xmark & \xmark & \xmark \\
\bottomrule
\end{tabular}
\end{table}

% ================= Method =================
\section{Method}
\label{sec:method}

\subsection{Problem setup}
Let $\mathbf{x}\in\mathcal{C}\subseteq\mathbb{R}^{d}$ be parameters constrained to a
feasible set $\mathcal{C}$ and $\mathbf{y}\in\mathbb{R}^{N}$ an observation from a
forward process $\mathbf{y}=\mathcal{F}(\mathbf{x})+\boldsymbol{\eta}$. The map is
many-to-one, so the target is the posterior $p(\mathbf{x}\mid\mathbf{y})$. We seek a
single model that (i)~predicts $\mathbf{y}$ from $\mathbf{x}$, (ii)~returns a point
estimate of $\mathbf{x}$, and (iii)~produces calibrated posterior samples, while
guaranteeing $\mathbf{x}\in\mathcal{C}$.

\subsection{A liquid-ODE flow and its latent reversibility}
The core is a continuous normalizing flow $F_\theta$ defined by integrating a learned
velocity field from $t{=}0$ to $t{=}1$, with the liquid form
\begin{equation}
g_\theta(\mathbf{v},t)= -\Big(\tfrac{1}{\boldsymbol{\tau}}+\mathbf{s}(\mathbf{v},t)\Big)\odot \mathbf{v}+\mathbf{s}(\mathbf{v},t)\odot \mathbf{A},\qquad \mathbf{s}(\mathbf{v},t)=\sigma\!\big(\mathrm{MLP}_\theta([\mathbf{v},t])\big),
\label{eq:liquid}
\end{equation}
where $\mathrm{MLP}_\theta$ uses a continuously differentiable activation. Fixing
$\mathbf{s}$ to be state-independent recovers a static neural ODE (an ablation).

% =====================================================================
% PRISM Phase 0 -- corrected theory (drop-in)
% Replaces old Proposition 1 (Sec. 3.2), the stability proposition (Sec. 3.6),
% and Appendix C proofs.
% Requires in preamble:
%   \usepackage{amsmath,amssymb,amsthm}
%   \newtheorem{proposition}{Proposition}
%   \newtheorem{assumption}{Assumption}
% =====================================================================

\begin{assumption}[Regularity of the velocity field]
\label{ass:reg}
The velocity field $g_\theta:\mathbb{R}^{d}\times[0,1]\to\mathbb{R}^{d}$ is
continuous in $t$, continuously differentiable in $\mathbf{v}$, and uniformly
Lipschitz in $\mathbf{v}$ on the relevant compact domain $\Omega$, with constant
$L_g=\sup_{t\in[0,1]}\mathrm{Lip}_{\mathbf{v}}\,g_\theta(\cdot,t)<\infty$. In the
implementation this is ensured by using a continuously differentiable activation
(SiLU) in $\mathrm{MLP}_\theta$; $\tanh$ or softplus satisfy the same conditions.
\end{assumption}

\begin{proposition}[Latent-flow reversibility]
\label{prop:inv}
Under Assumption~\ref{ass:reg}, the time-$1$ flow map
$F_\theta:\mathbf{v}(0)\mapsto\mathbf{v}(1)$ is a $C^{1}$ diffeomorphism onto its
image, and its inverse is obtained by integrating the same field in reverse
time; hence $F_\theta^{-1}\circ F_\theta=\mathrm{id}$ in exact arithmetic.
Moreover $\det\nabla F_\theta(\mathbf{v}(0))
=\exp\!\big(\int_{0}^{1}\mathrm{tr}\,\partial_{\mathbf{v}}
g_\theta(\mathbf{v}(t),t)\,dt\big)>0$, so $F_\theta$ is orientation-preserving
and everywhere invertible. This concerns the \emph{latent} transport only; it
does not by itself imply consistency between the auxiliary physical forward
surrogate $g_\psi$ and the inverse path, which we measure empirically.
\end{proposition}

\begin{proposition}[Numerical reconstruction error]
\label{prop:num}
Let $\widehat{F}_\theta$ denote integration of the flow with a stable one-step
method of order $p$ and step size $h$. Under Assumption~\ref{ass:reg}, the
forward--reverse reconstruction error satisfies
$\lVert \widehat{F}_\theta^{-1}(\widehat{F}_\theta(\tilde{\mathbf{x}}))
-\tilde{\mathbf{x}}\rVert=\mathcal{O}(h^{p})$, up to floating-point roundoff. We
make no claim of exact cancellation between the forward and reverse passes; the
error is governed by the global discretization error of the chosen integrator.
\end{proposition}

\begin{proposition}[Stability of the amortized inverse]
\label{prop:stab}
Let $h(\mathbf{y})=c\big(F_\theta^{-1}(\boldsymbol{\mu}_\phi(\mathbf{y}))\big)$.
Under Assumption~\ref{ass:reg}, a Gr\"onwall estimate gives
$\mathrm{Lip}(F_\theta^{-1})\le e^{L_g}$, hence
\begin{equation}
\mathrm{Lip}(h)\ \le\ \mathrm{Lip}(c)\,e^{L_g}\,\mathrm{Lip}(\boldsymbol{\mu}_\phi).
\end{equation}
This is a \emph{global} upper bound and may be loose; we therefore also report
the empirical local sensitivity
$\max_i\lVert\nabla h(\mathbf{y}_i)\rVert_2$ over held-out observations
(Table~\ref{tab:stability}) and distinguish it from the theoretical constant.
\end{proposition}

% ---------------------------------------------------------------------
% Appendix proofs
% ---------------------------------------------------------------------

\paragraph{Proof of Proposition~\ref{prop:inv}.}
By Assumption~\ref{ass:reg} and the Picard--Lindel\"of theorem, the initial value
problem $\dot{\mathbf{v}}(t)=g_\theta(\mathbf{v}(t),t)$,
$\mathbf{v}(0)=\tilde{\mathbf{x}}$, has a unique solution on $[0,1]$ depending
continuously on the initial condition. Since $g_\theta$ is $C^{1}$ in
$\mathbf{v}$, the variational equation
$\dot{\mathbf{J}}(t)=\partial_{\mathbf{v}}g_\theta(\mathbf{v}(t),t)\,\mathbf{J}(t)$,
$\mathbf{J}(0)=\mathbf{I}$, has solution
$\mathbf{J}(t)=\nabla F_\theta^{(t)}$ with
$\det\mathbf{J}(t)=\exp\!\big(\int_0^t\mathrm{tr}\,\partial_{\mathbf{v}}
g_\theta\,ds\big)>0$ by Liouville's formula. Hence $\nabla F_\theta$ is
nonsingular everywhere and, by the inverse function theorem, $F_\theta$ is a
$C^{1}$ diffeomorphism whose inverse is the time-reversed flow. \hfill$\square$

\paragraph{Proof of Proposition~\ref{prop:num}.}
Under Assumption~\ref{ass:reg}, a stable order-$p$ one-step method has global
error $\mathcal{O}(h^{p})$ on the fixed interval $[0,1]$. The forward and reverse
integrations compose two maps, each within $\mathcal{O}(h^{p})$ of the exact
(mutually inverse) flows; the triangle inequality gives a reconstruction error of
$\mathcal{O}(h^{p})$, up to roundoff. No symmetry of the integrator is assumed.
\hfill$\square$

\paragraph{Proof of Proposition~\ref{prop:stab}.}
For initial states $\mathbf{a},\mathbf{b}$,
$\frac{d}{dt}\lVert\Phi_t(\mathbf{a})-\Phi_t(\mathbf{b})\rVert
\le L_g\lVert\Phi_t(\mathbf{a})-\Phi_t(\mathbf{b})\rVert$, so Gr\"onwall's
inequality gives
$\lVert\Phi_1(\mathbf{a})-\Phi_1(\mathbf{b})\rVert
\le e^{L_g}\lVert\mathbf{a}-\mathbf{b}\rVert$, and the same bounds the reverse
flow. Sub-multiplicativity of Lipschitz constants over the composition
$c\circ F_\theta^{-1}\circ\boldsymbol{\mu}_\phi$ yields the stated bound.
\hfill$\square$


Empirically, the latent cycle error tracks the integration floor and the fitted
convergence order matches the integrator, while the \emph{physical} cycle
$\lVert h(g_\psi(\mathbf{x}))-\mathbf{x}\rVert$ is reported separately as an
empirical quantity (Table~\ref{tab:cycle}, Figure~\ref{fig:reversibility}).

\subsection{Hard constraints by construction}
% =====================================================================
% PRISM Phase 0 -- corrected constraint transforms (drop-in)
% Replaces Section 3.3 body and adds the transforms table + physical-space NLL.
% Requires: \usepackage{amsmath,amssymb,booktabs}
% =====================================================================

% ---- Section 3.3 body ----
We parameterize constrained parameters through an explicit, differentiable
support transform $c:\mathbb{R}^{k}\to\mathcal{C}$ whose domain dimension $k$
matches the \emph{intrinsic} dimension of the feasible set $\mathcal{C}$, with
closed-form inverse $c^{-1}$ and Jacobian. PRISM operates on the unconstrained
coordinate $\tilde{\mathbf{x}}=c^{-1}(\mathbf{x})$ and emits
$\mathbf{x}=c(\tilde{\mathbf{x}})$; since $c$ maps onto $\mathcal{C}$, every
decoded sample---including posterior tail draws---is feasible by construction.
Table~\ref{tab:transforms} lists the transforms. For positive parameters we use
the softplus map and for box-constrained parameters a scaled logistic map, both
smooth bijections onto the open feasible set. Simplex-valued parameters (not
required by our benchmarks) are supported through a dimension-matched
stick-breaking bijection $\mathbb{R}^{d-1}\to\mathrm{int}\,\Delta^{d-1}$; we
avoid the softmax, which is not injective (adding a constant to every logit
leaves the output unchanged) and maps a $d$-dimensional input onto a set of
intrinsic dimension $d-1$.

\paragraph{Density in physical coordinates.}
Because the conditional density is defined on the unconstrained coordinate, the
physical-space log-density includes the transform Jacobian,
\begin{equation}
\log p_\theta(\mathbf{x}\mid\mathbf{y})
=\log p_\theta\big(c^{-1}(\mathbf{x})\mid\mathbf{y}\big)
+\log\big|\det \mathbf{J}_{c^{-1}}(\mathbf{x})\big|,
\label{eq:phys-nll}
\end{equation}
and for elementwise transforms
$\log|\det\mathbf{J}_{c^{-1}}(\mathbf{x})|
=\sum_i\log|(c^{-1})'(x_i)|=-\sum_i\log|c'(u_i)|$ with $u_i=c^{-1}(x_i)$. All
reported negative log-likelihoods are evaluated in physical coordinates with this
term included; we verify Eq.~\eqref{eq:phys-nll} against a one-dimensional
numerical normalization check for the softplus and logistic maps.

% ---- Transforms table (appendix or Sec. 3.3) ----
\begin{table}[t]
\centering
\caption{Support transforms used by PRISM. Each is a smooth bijection from the
unconstrained coordinate onto the (open) feasible set, with closed-form inverse
and Jacobian; the domain dimension matches the intrinsic dimension of the
feasible set. Here $\sigma$ is the logistic sigmoid and $z=(x-a)/(b-a)$.}
\label{tab:transforms}
\begin{tabular}{lllll}
\toprule
Constraint & $c(u)$ & Codomain & $c^{-1}(x)$ & $\log|c'(u)|$ \\
\midrule
Positive & $\mathrm{softplus}(u)$ & $(0,\infty)$ & $\log(e^{x}-1)$ & $\log\!\big(1-e^{-x}\big)$ \\
Box $[a,b]$ & $a+(b-a)\,\sigma(u)$ & $(a,b)$ & $\sigma^{-1}\!\big(z\big)$ & $\log(b-a)+\log z+\log(1-z)$ \\
Simplex$^{\dagger}$ & stick-breaking & $\mathrm{int}\,\Delta^{d-1}$ & log-ratio & closed form \\
\bottomrule
\end{tabular}
\\[3pt]
{\footnotesize $^{\dagger}$Not used by our benchmarks; listed to show
extensibility with a dimension-matched bijection. The softmax is excluded as it
is non-injective and dimension-mismatched.}
\end{table}


\subsection{Calibrated posterior via an amortized conditional base}
A continuous flow induces a density through the instantaneous change of variables.
We condition on $\mathbf{y}$ in the \emph{base distribution}: an amortization
network maps $\mathbf{y}$ to a Gaussian base, giving
\begin{equation}
\log p_\theta(\tilde{\mathbf{x}}\mid\mathbf{y})=\log\mathcal{N}\!\big(F_\theta(\tilde{\mathbf{x}});\,\boldsymbol{\mu}_\phi(\mathbf{y}),\,\mathrm{diag}\,\boldsymbol{\sigma}_\phi(\mathbf{y})^2\big)+\log\big|\det\nabla_{\tilde{\mathbf{x}}}F_\theta(\tilde{\mathbf{x}})\big|.
\label{eq:cond}
\end{equation}
We show empirically (Table~\ref{tab:conditioning}) that conditioning only through
the velocity field can, in our setup, ignore the observation; conditioning the base
yields informative, y-dependent posteriors. Sampling is amortized: draw
$\mathbf{z}\sim\mathcal{N}(\boldsymbol{\mu}_\phi(\mathbf{y}),\cdot)$ and return
$\mathbf{x}=c(F_\theta^{-1}(\mathbf{z}))$.

\subsection{Forward surrogate and training}
A lightweight head $g_\psi$ provides the forward map. Training minimizes the
physical-space negative log-likelihood (NLL), including the transform Jacobian of
Eq.~\eqref{eq:phys-nll}, plus a forward term, with Adam.

% ================= Experiments =================
\section{Experiments}
\label{sec:experiments}
We evaluate on eight procedurally generated benchmarks across PDE operator
inversion (Darcy, Burgers, Helmholtz), simulation-based inference (Gaussian-Linear,
SLCP, Two Moons), and dynamical systems (Oscillator, Lotka-Volterra). SLCP has a
single dominant mode; Two Moons is the disconnected bimodal case. Baselines are
cINN, NPE with a continuous normalizing flow (NPE-CNF), amortized variational
inference (VI), an MCMC reference, and the forward operators FNO, DeepONet, and a
PINN. All learned-method tables report mean $\pm$ standard deviation over five
seeds, and every value is read from a single result manifest. Full configuration
is in Appendix~\ref{app:repro}.

\subsection{Reversibility (latent, guaranteed; physical, measured)}
\input{paper_assets/T_cycle_phase2.tex}
Figure~\ref{fig:reversibility} shows the latent cycle error following the expected
$\mathcal{O}(h^{p})$ convergence to an integration floor, while the physical cycle
is reported separately and never presented as exact.

\subsection{Feasibility: a fair comparison}
\input{paper_assets/T_constraints_fair.tex}
When every method receives the same support transform, the violation rate is
$\approx 0$ across the board; the comparison therefore turns on calibration,
physical NLL, inverse error, and latency rather than on feasibility.

\subsection{Stability}
\input{paper_assets/T_stability.tex}
The global Gr\"onwall bound is loose; the empirical local sensitivity is far
smaller and governs practical robustness (Figure~\ref{fig:stability}).

\subsection{Ablations}
\input{paper_assets/T_liquid_vs_static.tex}
\input{paper_assets/T_conditioning.tex}
\input{paper_assets/T_projection.tex}
Liquid dynamics help most on the harder posterior geometries rather than
uniformly; the conditioning ablation shows the dynamics-only variant ignoring the
observation; and the projection's value is visible on the box-constrained tasks.

\subsection{Accuracy}
\input{paper_assets/T_main_pde.tex}
On the forward map a dedicated operator is most accurate; PRISM's auxiliary head is
competitive but not state-of-the-art, and it leads on the inverse RMSE this work
targets (Figure~\ref{fig:fields}).

\subsection{Calibration across domains}
\input{paper_assets/T_calibration_ext.tex}
PRISM is competitive with the strongest learned baselines and has strong coverage
and C2ST on several tasks, while VI attains the lowest expected coverage error on
the Gaussian tasks (Figure~\ref{fig:calibration}).

\subsection{Two Moons: joint geometry}
\input{paper_assets/T_joint_geometry.tex}
Marginal coverage is acceptable, yet the joint discrepancy is large and the mode
balance low: a single connected base blurs the two crescents
(Figure~\ref{fig:posteriors}). We do not present marginal calibration as evidence
of a correct joint posterior.

\subsection{Efficiency and generality}
\input{paper_assets/T_efficiency.tex}
\input{paper_assets/T_generality_raw.tex}
PRISM trades higher per-sample latency for continuous-time reversibility and hard
constraints; the amortization break-even falls sharply as parameter dimension grows
(Figure~\ref{fig:pareto}). Generality is reported as raw metrics with a
posterior-appropriate primary metric per task, not a binary pass/fail.

% ================= Limitations & Conclusion =================
\section{Limitations and Conclusion}
\label{sec:limitations}
A single continuous flow with a unimodal base cannot split disconnected posterior
modes; on Two Moons PRISM recovers a calibrated but mode-blurred posterior. Trained
on a parameter sub-range, amortized inference does not extrapolate out of
distribution. Each posterior sample requires an ODE solve, so per-sample latency
exceeds that of discrete flows, though amortization wins in total wall-clock beyond
a break-even that shrinks with dimension. The forward surrogate is competitive but
not state-of-the-art. Within these bounds, PRISM shows that reversible latent
transport, feasibility by construction, and amortized uncertainty-aware inference
can coexist in one model, with guarantees where they are provable and measurement
where they are not.

% =====================================================================
% PRISM Phase 6 -- reproducibility statement + appendix (drop-in)
% Place the statement before References; place the appendix section in the
% appendix. Fills guide item 22.
% =====================================================================

% ---- Reproducibility Statement (before \bibliography) ----
\subsubsection*{Reproducibility Statement}
All benchmarks are procedurally generated, so no external data are required. The
method is fully specified in Section~\ref{sec:method} (the velocity field,
support transforms, conditional base, training objective, and the physical-space
likelihood) and the theory in Appendix~\ref{app:proofs} states every regularity
assumption. Appendix~\ref{app:repro} lists dataset sizes and splits, architecture
widths and activations, optimizer and solver settings, the Jacobian estimator,
posterior sample counts, the five random seeds, the hyperparameter-search budget,
and hardware. Every number in the paper is written to a single result manifest
and read back by the table and figure generators, so a value cannot appear twice
with two figures; the entire set of tables and figures regenerates from the
frozen manifest with one command, and all experiments re-run from scratch with
one script.

% ---- Appendix: full reproducibility details ----
\section{Reproducibility details}
\label{app:repro}

\paragraph{Data.} Each task is generated on the fly from a fixed seed. Training
uses $N_{\text{train}}$ parameter--observation pairs, validation $N_{\text{val}}$,
and test $N_{\text{test}}$, drawn i.i.d.\ from the task prior and forward process
(Table~\ref{tab:repro-data}). Because generation is procedural, simulation-based
calibration draws fresh $(\theta, y)$ pairs at evaluation time.

\paragraph{Architecture.} The liquid velocity field $g_\theta$ is a two-hidden-layer
multilayer perceptron with SiLU activations and width $W$; the amortization
network $(\boldsymbol{\mu}_\phi, \boldsymbol{\sigma}_\phi)$ and the forward head
share this width. The flow is integrated with a fixed-step Runge--Kutta method of
order $p$ using $K$ steps; the Jacobian trace uses an exact evaluation in the
low-dimensional regime and a Hutchinson estimator otherwise. Constraints use the
support transforms of Table~\ref{tab:transforms}.

\paragraph{Optimization.} All methods are trained with Adam at learning rate
$\eta$, batch size $B$, for $E$ epochs, with gradient clipping at $5.0$ and a
cosine schedule. Baselines receive method-specific tuning with an equal search
budget; the NPE-CNF baseline is validated on the analytic linear-Gaussian task
before use.

\paragraph{Evaluation.} Posterior metrics use $S$ samples per observation.
Simulation-based calibration uses $N_{\text{SBC}}$ observations. All learned-method
tables report mean $\pm$ standard deviation over seeds $\{0,1,2,3,4\}$.

\paragraph{Hardware and runtime.} Experiments run on \texttt{<GPU/CPU model>};
the full five-seed sweep completes in \texttt{<hh:mm>}. Exact commands to
regenerate each table and figure are listed in the repository \texttt{README}.

\begin{table}[h]
\centering
\caption{Dataset sizes and key hyperparameters (fill from your configuration).}
\label{tab:repro-data}
\begin{tabular}{lcccccccc}
\toprule
Task & $N_{\text{train}}$ & $N_{\text{val}}$ & $N_{\text{test}}$ & $W$ & $K$ & $\eta$ & $B$ & $E$ \\
\midrule
Darcy          & & & & & & & & \\
Burgers        & & & & & & & & \\
Helmholtz      & & & & & & & & \\
Gaussian-Linear & & & & & & & & \\
SLCP           & & & & & & & & \\
Two Moons      & & & & & & & & \\
Oscillator     & & & & & & & & \\
Lotka-Volterra & & & & & & & & \\
\bottomrule
\end{tabular}
\end{table}


\bibliography{references}
\bibliographystyle{iclr2026_conference}

\appendix
\section{Additional figures}
\begin{figure}[h]\centering\includegraphics[width=\linewidth]{paper_assets/F_reversibility.pdf}
\caption{Latent (guaranteed) and physical (empirical) cycle error versus solver
budget. The latent error follows the expected convergence order down to an
integration floor; the physical cycle is a separate, measured quantity.}
\label{fig:reversibility}\end{figure}
\begin{figure}[h]\centering\includegraphics[width=0.7\linewidth]{paper_assets/F_stability.pdf}
\caption{Empirical local sensitivity (mean and max) versus perturbation, with the
loose global bound for context. We do not equate the empirical max with the true
global Lipschitz constant.}
\label{fig:stability}\end{figure}
\begin{figure}[h]\centering\includegraphics[width=\linewidth]{paper_assets/F_calibration_ext.pdf}
\caption{Calibration across five domains: SBC rank histograms (near-uniform
indicates calibration) and reliability curves tracking the diagonal.}
\label{fig:calibration}\end{figure}
\begin{figure}[h]\centering\includegraphics[width=\linewidth]{paper_assets/F_posteriors_tm.pdf}
\caption{Two Moons posteriors: reference versus PRISM. A single connected base
blurs the modes; the mixture base, where used, recovers them.}
\label{fig:posteriors}\end{figure}
\begin{figure}[h]\centering\includegraphics[width=0.5\linewidth]{paper_assets/F_fields.pdf}
\caption{Darcy field: ground truth versus PRISM reconstruction (shared color
scale).}
\label{fig:fields}\end{figure}
\begin{figure}[h]\centering\includegraphics[width=0.55\linewidth]{paper_assets/F_pareto.pdf}
\caption{Accuracy versus latency, drawn from the same records as the efficiency
table.}
\label{fig:pareto}\end{figure}

\end{document}
